% Auto-generated by scripts/exp_fold_gauge_anomaly_ldp_rate.py \begin{proposition}[规范差大偏差率函数接口(Legendre 变换 + 端点闭式)]\label{prop:fold-gauge-anomaly-ldp} Under the uniform baseline $\omega\sim\mathrm{Unif}(\Omega_{m+1})$, let $G_m=\sum_{t=1}^{m}g_t$, 其中 $g_t=\mathbf 1\{X_t\neq Y_t\}$ denote the gauge discrepancy逐位失配指示.then对任意 $a\in[0,1]$, 比例 $G_m/m$ 满足大偏差原理, 其速率函数可由压力函数做 Legendre--Fenchel 变换给出:  \[ I(a)=\sup_{\theta\in\mathbb R}\bigl(a\theta-P_G(\theta)\bigr),\qquad P_G(\theta)=\log\rho(A_\theta), \] 其中 $A_\theta$ 为命题 \ref{prop:fold-gauge-anomaly-pressure} $4\times 4$ 加权邻接矩阵.并且端点稀有事件指数具有完全闭式:  \[ \boxed{\ I(0)=\log(2/\varphi),\qquad I(1)=\log 2\ }. \] \end{proposition} \begin{proof}[Proof sketch(auditable)] 由于 $g_t$ 是有限状态边移位上可加观测量, G\"artner--Ellis 定理给出 $I$ 为 $P_G$ Legendre--Fenchel 变换.端点闭式来自极限扭曲: 当 $\theta\to-\infty$(即 $u=e^\theta\to 0$)时, 失配边权趋于 $0$, 谱半径退化为“全匹配子图” Perron 根, 给出 $P_G(-\infty)=\log(\varphi/2)$, 故 $I(0)=-P_G(-\infty)=\log(2/\varphi)$; 当 $\theta\to+\infty$ 时, 谱半径由全失配三环支配, 满足 $P_G(\theta)=\theta-\log 2+o(1)$, 故 $I(1)=\log 2$. \end{proof} 